% GENERATED FILE. Edit canonical lessons, then run npm run generate:books.
% canonical-source-hash: fnv1a64:20bfb2f33f5488cc
% canonical-lessons: GE-C90-schwer, GE-C90-billig, GE-C90-hoch, GE-C90-niedrig, GE-C90-leer

\chapter{Heavy, Cheap, High, Low, Empty}
\label{ch:ge-heavy-cheap-high-low-empty}
\hlchaptermodality{90}

\begin{chapteropening}
\textbf{By the end of this chapter:} \emph{I can say heavy, cheap, high, tall, low and empty.}

Complete the last lesson of chapter 90: I can say heavy, cheap, high, tall, low and empty.
\end{chapteropening}

\section[schwer]{schwer — heavy; difficult}

\label{lesson:GE-C90-schwer}



\begin{quote}
Before the new one: say the German for slow, then the German for light; easy.
\end{quote}

\subsection*{You'll want to know: schwer}
\textbf{schwer} — \textquotedblleft{}heavy; difficult\textquotedblright{}.

\begin{cousinweb}
\textbf{schwer} $\leftarrow$ Old High German \emph{swār} $\leftarrow$ Proto-Germanic \emph{*swēraz}, \textquotedblleft{}heavy\textquotedblright{} — inherited. (EtymDB 2.1)
\end{cousinweb}

\begin{grammarlens}[title={What you've built}]
One more word for this chapter.
\end{grammarlens}

\subsection*{Your turn}
\begin{itemize}
\raggedright
  \item \emph{Say it:} \emph{schwer}
  \item \emph{Say it:} \emph{schwer}, once more
  \item \emph{Say it:} say \emph{schwer}
  \item \emph{Recall:} say the German for slow, then the German for light; easy, then say \emph{schwer} again
\end{itemize}

\begin{tcolorbox}[breakable,colback=teal!4,colframe=teal!35!black,title={Before you move on}]
What does schwer mean? (\textquotedblleft{}Heavy; difficult\textquotedblright{}.) Say it once more.
\end{tcolorbox}
\section[billig]{billig — cheap}

\label{lesson:GE-C90-billig}



\begin{quote}
Before the new one: say the German for light; easy, then the German for heavy; difficult.
\end{quote}

\subsection*{You'll want to know: billig}
\textbf{billig} — \textquotedblleft{}cheap\textquotedblright{}.

\begin{cousinweb}
\textbf{billig} $\leftarrow$ Old High German \emph{billīh}, \textquotedblleft{}appropriate, fair\textquotedblright{} (EtymDB 2.1) — inherited. A fair price slid into a cheap one.
\end{cousinweb}

\begin{grammarlens}[title={What you've built}]
One more word for this chapter.
\end{grammarlens}

\subsection*{Your turn}
\begin{itemize}
\raggedright
  \item \emph{Say it:} \emph{billig}
  \item \emph{Say it:} \emph{billig}, once more
  \item \emph{Say it:} say \emph{billig}
  \item \emph{Recall:} say the German for light; easy, then the German for heavy; difficult, then say \emph{billig} again
\end{itemize}

\begin{tcolorbox}[breakable,colback=teal!4,colframe=teal!35!black,title={Before you move on}]
What does billig mean? (\textquotedblleft{}Cheap\textquotedblright{}.) Say it once more.
\end{tcolorbox}
\section[hoch]{hoch — high, tall}

\label{lesson:GE-C90-hoch}



\begin{quote}
Before the new one: say the German for heavy; difficult, then the German for cheap.
\end{quote}

\subsection*{You'll want to know: hoch}
\textbf{hoch} — \textquotedblleft{}high, tall\textquotedblright{}.

\begin{cousinweb}
\textbf{hoch} $\leftarrow$ Old High German \emph{hōh}, \textquotedblleft{}high, tall\textquotedblright{} $\leftarrow$ Proto-Germanic \emph{*hauhaz} — inherited, the same word as English \emph{high}. (EtymDB 2.1)
\end{cousinweb}

\begin{grammarlens}[title={What you've built}]
One more word for this chapter.
\end{grammarlens}

\subsection*{Your turn}
\begin{itemize}
\raggedright
  \item \emph{Say it:} \emph{hoch}
  \item \emph{Say it:} \emph{hoch}, once more
  \item \emph{Say it:} say \emph{hoch}
  \item \emph{Recall:} say the German for heavy; difficult, then the German for cheap, then say \emph{hoch} again
\end{itemize}

\begin{tcolorbox}[breakable,colback=teal!4,colframe=teal!35!black,title={Before you move on}]
What does hoch mean? (\textquotedblleft{}High, tall\textquotedblright{}.) Say it once more.
\end{tcolorbox}
\section[niedrig]{niedrig — low}

\label{lesson:GE-C90-niedrig}



\begin{quote}
Before the new one: say the German for cheap, then the German for high.
\end{quote}

\subsection*{You'll want to know: niedrig}
\textbf{niedrig} — \textquotedblleft{}low\textquotedblright{}.

\begin{cousinweb}
\textbf{niedrig} is the word for \textquotedblleft{}down\textquotedblright{} + \emph{-ig}. That word $\leftarrow$ Old High German \emph{nidar} $\leftarrow$ Proto-Germanic \emph{*niþer} (EtymDB 2.1); English \emph{nether}, \textquotedblleft{}lower\textquotedblright{}, is a cousin.
\end{cousinweb}

\begin{grammarlens}[title={What you've built}]
One more word for this chapter.
\end{grammarlens}

\subsection*{Your turn}
\begin{itemize}
\raggedright
  \item \emph{Say it:} \emph{niedrig}
  \item \emph{Say it:} \emph{niedrig}, once more
  \item \emph{Say it:} say \emph{niedrig}
  \item \emph{Recall:} say the German for cheap, then the German for high, then say \emph{niedrig} again
\end{itemize}

\begin{tcolorbox}[breakable,colback=teal!4,colframe=teal!35!black,title={Before you move on}]
What does niedrig mean? (\textquotedblleft{}Low\textquotedblright{}.) Say it once more.
\end{tcolorbox}
\section[leer]{leer — empty}

\label{lesson:GE-C90-leer}



\begin{quote}
Before the new one: say the German for high, then the German for low.
\end{quote}

\subsection*{You'll want to know: leer}
\textbf{leer} — \textquotedblleft{}empty\textquotedblright{}.

\begin{cousinweb}
\textbf{leer} $\leftarrow$ Middle High German \emph{lære}, \textquotedblleft{}empty\textquotedblright{} (EtymDB 2.1) — inherited, from Old High German \emph{lāri}.
\end{cousinweb}

\begin{grammarlens}[title={What you've built}]
One more word for this chapter.
\end{grammarlens}

\subsection*{Your turn}
\begin{itemize}
\raggedright
  \item \emph{Say it:} \emph{leer}
  \item \emph{Say it:} \emph{leer}, once more
  \item \emph{Say it:} say \emph{leer}
  \item \emph{Recall:} say the German for high, then the German for low, then say \emph{leer} again
\end{itemize}

\begin{tcolorbox}[breakable,colback=teal!4,colframe=teal!35!black,title={Before you move on}]
What does leer mean? (\textquotedblleft{}Empty\textquotedblright{}.) Say it once more.
\end{tcolorbox}
